\paragraph{Goal Completion}
Since social goals can be subjective, we judge its completion from three different aspects, as shown in Figure~\ref{fig:fm}:
(1) \textbf{Self}: We ask the agent whether it has achieved its goals after interacting with others.
(2) \textbf{Other}: We ask other agents in the scenario whether the target agent has achieved his or her specific goals.
(3) \textbf{External}: We prompt third-party models with the chat history and ask if an agent has achieved his/her own social goals.
We ask the interviewees to respond with \textit{yes} (goal completed) or \textit{no}.
We take the average across all goals of a character to measure the agent's overall goal completion level.

\paragraph{Implicit Reasoning}
As mentioned in Section~\ref{sec:scenario_construction}, each character's private information corresponds to a multiple-choice evaluation question.
To evaluate an agent's information reasoning ability, we present it with questions related to the private information of other agents within the scenario.
We then calculate the average accuracy (Acc) of the current agent on these questions to determine the agent's score in information reasoning.

\paragraph{Profile Sensitivity} 
After character enrichment, each template generates multiple scenarios. By incorporating diverse characters, we not only enrich the scenarios but also gain insights into the stability of social intelligence when simulating different roles.
Thus, we propose \textbf{profile sensitivity index (PSI)}. We compute the standard deviation (std) of goal/information metrics of the scenario sharing the same template, and the average std across all templates is calculated as PSI.
A lower PSI indicates that social intelligence is more stable.